%% Springer Nature (sn-jnl) version for Empirical Software Engineering.
%% GENERATED by experiments/build_ese.py from manuscript/latex/*.tex --
%% edit the sources there (or this template), not main.tex.
\documentclass[pdflatex,sn-basic]{sn-jnl}

\usepackage{graphicx}%
\usepackage{multirow}%
\usepackage{amsmath,amssymb,amsfonts}%
\usepackage{amsthm}%
\usepackage{mathrsfs}%
\usepackage{xcolor}%
\usepackage{textcomp}%
\usepackage{manyfoot}%
\usepackage{booktabs}%
\usepackage{algorithm}%
\usepackage{algorithmicx}%
\usepackage{algpseudocode}%
\usepackage{tikz}%
\usetikzlibrary{positioning,arrows.meta,fit,calc,backgrounds}%
\usepackage{url}%

\theoremstyle{thmstyleone}%
\newtheorem{proposition}{Proposition}%
\raggedbottom

%%NUMBERS%%

\begin{document}

\title[Node descriptions for cross-project vulnerability detection]{What Should a Graph Node Say? A Controlled Study of Node Descriptions for Cross-Project Vulnerability Detection}

\author*[1]{\fnm{Quang-Vinh} \sur{Dang}}\email{vinh.dq4@buv.edu.vn}

\affil*[1]{\orgname{British University Vietnam}, \orgaddress{\city{Hung Yen}, \country{Vietnam}}}

\abstract{\textbf{Context:} Learned vulnerability detectors often fail on software projects that were not in their training data. \textbf{Objective:} We ask how much of that failure depends on what each node of the program graph is told about the code. \textbf{Method:} We change only the node description: identifier text, the same text with identifiers renamed, the name of the grammar rule that produced the node (the grammar kind), and a hand-built label for the operation a node performs. Detectors are trained on uniform random samples of two public corpora (BigVul and DiverseVul, about 10{,}000 functions each) and tested on held-out whole projects, over ten matched folds, with a cluster bootstrap over projects as the primary inference. Performance is measured by AUC, the chance that a vulnerable function is ranked above a benign one (0.5 is random). \textbf{Results:} Text read from scratch ranks functions close to random (AUC \mBvGgnnLexAuc{} and \mDvGgnnLexAuc{} with a gated graph network). Renaming identifiers or using grammar kinds raises AUC by roughly \numdeltaMin{} to \numdeltaMax{} on both corpora and under every architecture tried, and the two work about equally well. The hand-built labels add nothing to the grammar kind; they are \numdeterminismMin{}--\numdeterminismMax{}\% predictable from it. Much of the gain disappears once function size is held fixed, and a rule that ranks functions only by length reaches AUC \numbvSizeAuc{} and \numdvSizeAuc{}, above every model. \textbf{Conclusions:} Studies of cross-project detectors should report function-size and vocabulary-free baselines; we release code, annotations and scripts.}

\keywords{Vulnerability detection, Cross-project evaluation, Program representation, Abstract syntax trees, Graph neural networks}

\maketitle

%%BODY%%

\section*{Declarations}

\begin{itemize}
\item \textbf{Funding:} This research did not receive any specific grant from funding agencies in the public, commercial, or not-for-profit sectors.
\item \textbf{Conflict of interest/Competing interests:} The author has no relevant financial or non-financial interests to disclose.
\item \textbf{Ethics approval and consent to participate:} Not applicable. The study uses publicly available source-code corpora and involves no human participants.
\item \textbf{Consent for publication:} Not applicable.
\item \textbf{Data availability:} The study uses the public corpora BigVul, DiverseVul and Devign, accessed through the public mirrors named in the text. The sampled function identifiers, the per-fold held-out predictions of every system, and the annotation sheets are in the repository below.
\item \textbf{Materials availability:} Not applicable.
\item \textbf{Code availability:} All code, annotation sheets and the scripts that regenerate every table, figure and number in this paper are available at \url{https://github.com/vinhqdang/vulmorph-fed}.
\item \textbf{Author contribution:} Single author.
\item \textbf{Use of large language models:} Language models were used as the two blind annotators and the adjudicator in the validation of Section~\ref{sec:rq_analysis}, as described there. Language-model assistants also helped with writing code and drafting text; the author reviewed and edited all content and takes full responsibility for it.
\end{itemize}

\bibliography{ref}

\end{document}
